\documentclass[12pt]{article}
\usepackage[utf8]{inputenc}
\usepackage[T1]{fontenc}
\usepackage{hyperref}
\title{Western Tamang Language Resources for K-12 Curriculum Development}
\author{Compiled by Claude Sonnet 5 (claude-sonnet-5) for Andr\'as Kornai}
\date{}
\begin{document}
\maketitle
\section{Introduction}
Western Tamang (ISO 639-3 tdg, Glottocode west2415) is a Sino-Tibetan language of Nepal. Tamang was reported at 5.1 percent of Nepal's population in the 2011 census. This report is based on one web search.
\section{Primary Educational Resources}
No mother-tongue textbook or primer was found in the search results.
\section{Linguistic and Reference Materials}
\begin{itemize}
\item \emph{Western Tamang Dictionary} (2017), Western Tamang--Nepali, SIL International and Mother Tongue Center Nepal. It uses a 31-letter Nepali-script alphabet and is available online and for download. Overview: \url{https://www.webonary.org/wtamang/overview/?lang=en}
\item Introduction to the dictionary: \url{https://www.webonary.work/wtamang/overview/introduction/?lang=en}
\item Alphabet: \url{https://www.webonary.org/wtamang/overview/alphabet/}
\item Download page (app and PDF): \url{https://www.webonary.org/wtamang/download/?lang=en}
\item Everitt, Taylor and Tamang (1972), a 40-page Tamang vocabulary, Glottolog reference record: \url{https://glottolog.org/resource/reference/id/26744}
\end{itemize}
\section{Religious and Community Resources}
Mother Tongue Center Nepal is named as co-publisher of the dictionary. No other community resources were found.
\section{Digital and Audio Resources}
The dictionary is available as a web edition and a downloadable app (see the download page above). No audio resources were found.
\section{Curriculum Development Context}
The language has a written orthography and a modern bilingual dictionary. The 2017 dictionary draws on a 2016 workshop with 29 speakers from six districts. No textbook or curriculum documents were found, despite the language being taught as a subject in some schools. Searching Nepal's Curriculum Development Centre and Mother Tongue Center Nepal materials would be the next step.
\end{document}